%% ===================================================================
\section{Método}\label{sec:metodo}

Esta revisión sigue el procedimiento de Kitchenham y Charters para revisiones sistemáticas en ingeniería de software \cite{kitchenham2007} y se reporta conforme a la declaración PRISMA 2020 \cite{page2021prisma}. El procedimiento se organiza en tres fases (planificación, conducción y reporte) cuyas actividades se resumen en la Tabla~\ref{tab:proceso}. Antes de la búsqueda se elaboró un protocolo que fijó las preguntas de investigación, la estrategia de búsqueda, los criterios de selección, la evaluación de la calidad y el proceso de extracción y síntesis. \pend{Indicar si el protocolo se registró y, en su caso, el repositorio, el identificador y las enmiendas.} Cada etapa se describe en los apartados siguientes.

\begin{table}[t]
\caption{Fases del proceso de revisión}\label{tab:proceso}
\begin{tabular*}{\textwidth}{@{\extracolsep\fill}ll}
\toprule
Fase & Actividades \\
\midrule
Planificación & Identificación de la necesidad de la revisión \\
 & Formulación de las preguntas de investigación \\
 & Elaboración del protocolo de revisión \\
Conducción & Búsqueda e identificación de estudios \\
 & Selección de estudios primarios \\
 & Evaluación de la calidad \\
 & Extracción de datos \\
 & Síntesis de datos \\
Reporte & Redacción y verificación del informe según PRISMA 2020 \\
\botrule
\end{tabular*}
\end{table}


\subsection{Preguntas de investigación}

El primer paso del protocolo fue formular las preguntas que orientan la revisión (Tabla~\ref{tab:rq}). Su alcance se delimitó con el marco PICO:
\begin{itemize}
\item \textbf{Población:} datos visuales protegidos mediante perturbación proactiva.
\item \textbf{Intervención:} contramedidas aplicadas a esos datos o a su entrenamiento.
\item \textbf{Comparación:} entrenamiento con datos protegidos sin la contramedida, complementado con el entrenamiento limpio en condiciones comparables.
\item \textbf{Resultado:} persistencia del efecto protector y recuperación del rendimiento.
\end{itemize}

\begin{table}[t]
\caption{Preguntas de investigación}\label{tab:rq}
\footnotesize
\begin{tabularx}{\textwidth}{@{}L{0.9cm}Y{1.0}@{}}
\toprule
ID & Pregunta de investigación \\
\midrule
RQ1 & ¿Qué estrategias y algoritmos de perturbación proactiva se han propuesto para dificultar el entrenamiento no autorizado de modelos discriminativos con datos visuales, cómo se caracterizan y en qué tareas, conjuntos de datos y arquitecturas se han evaluado? \\
RQ2 & ¿Qué contramedidas se han evaluado frente a estas protecciones y bajo qué conocimientos, capacidades, acceso y recursos del adversario? \\
RQ3 & ¿Qué grado de protección persiste y qué recuperación del rendimiento logra el adversario tras aplicar las contramedidas evaluadas, en los distintos dominios y configuraciones experimentales? \\
RQ4 & ¿Qué referencias experimentales y métricas se utilizan para evaluar la robustez y la reversión adversaria, y en qué condiciones permiten comparar los resultados entre estudios y dominios? \\
\botrule
\end{tabularx}
\end{table}


Cada pregunta cumple una función distinta:
\begin{itemize}
\item \textbf{RQ1} caracteriza las protecciones propuestas: cómo se generan, qué restricciones imponen, en qué modalidades y tareas se aplican y con qué conjuntos de datos y arquitecturas se evalúan. Esta caracterización es necesaria para interpretar las respuestas posteriores, porque la robustez de una protección depende de su mecanismo.
\item \textbf{RQ2} identifica las contramedidas evaluadas y el modelo de adversario bajo el que se aplican. Una contramedida solo es relevante si sus supuestos sobre el conocimiento, el acceso y los recursos del adversario son verosímiles para el escenario de protección.
\item \textbf{RQ3} cuantifica cuánto rendimiento recupera el adversario y cuánto efecto protector persiste. Es la pregunta central de la revisión.
\item \textbf{RQ4} examina las referencias experimentales y las métricas con las que se mide la robustez, para determinar en qué condiciones los resultados de estudios distintos pueden compararse.
\end{itemize}

\subsection{Estrategia de búsqueda}

\subsubsection{Fuentes de información}

Se consultaron las cuatro bases de datos de la Tabla~\ref{tab:bases}. Se eligieron por su cobertura de las revistas y actas de congresos de aprendizaje automático, visión por computadora y seguridad, y porque permiten búsquedas booleanas sobre los metadatos y la exportación de los registros. Como fuente complementaria se incorporó una colección de artículos conocidos por el equipo, consultada en la plataforma Undermind, que se reporta como otros métodos.

\begin{table}[t]
\caption{Fuentes de información}\label{tab:bases}
\footnotesize
\begin{tabular*}{\textwidth}{@{\extracolsep\fill}lll@{}}
\toprule
Base de datos & Campos de búsqueda & URL \\
\midrule
Scopus & Título, resumen y palabras clave & \url{https://www.scopus.com/} \\
Web of Science & Tema (título, resumen y palabras clave) & \url{https://www.webofscience.com/} \\
IEEE Xplore & Todos los metadatos & \url{https://ieeexplore.ieee.org/} \\
SpringerLink & Búsqueda general de la plataforma & \url{https://link.springer.com/} \\
\botrule
\end{tabular*}
\end{table}


\subsubsection{Cadena de búsqueda}

La cadena combinó cinco conceptos mediante el operador OR (Tabla~\ref{tab:cadena}):
\begin{itemize}
\item la terminología propia de los datos no aprendibles;
\item los ataques de disponibilidad, restringidos a contextos de aprendizaje con imágenes;
\item las perturbaciones protectoras o proactivas;
\item el envenenamiento indiscriminado y contrastivo;
\item el camuflaje frente al reconocimiento facial.
\end{itemize}
Los ataques de disponibilidad se acotaron con términos de imagen y aprendizaje porque, sin esa restricción, el término recupera estudios de disponibilidad de servicios de red. La cadena se adaptó a la sintaxis de cada base y se aplicaron los filtros de periodo (2020--2026) y de tipo documental (artículo de revista y ponencia de congreso). En Web of Science, por ejemplo, la consulta fue:

\begin{quote}
\footnotesize
TS=(``unlearnable example*'' OR ``unlearnable data'' OR ``unlearnable dataset*'' OR ``unlearnable sample*'' OR ``unlearnability'' OR ``error-minimizing noise'' OR ``unexploitable data'' OR ``robust unlearnable example*'' OR ((``availability attack*'' OR ``availability poison*'') AND (image* OR ``training data'' OR ``deep learning'' OR ``machine learning'')) OR ``data-protecting perturbation*'' OR ``protective perturbation*'' OR ``proactive perturbation*'' OR ``indiscriminate poison*'' OR ``contrastive poison*'' OR (cloak* AND (``face recognition'' OR ``facial recognition''))) AND PY=(2020-2026)
\end{quote}

Las búsquedas se ejecutaron en octubre de 2026. Las cadenas completas de cada base se incluyen como material suplementario.

\begin{table}[t]
\caption{Cadena de búsqueda}\label{tab:cadena}
\footnotesize
\begin{tabularx}{\textwidth}{@{}L{3.2cm}Y{1.0}@{}}
\toprule
Concepto & Términos de búsqueda \\
\midrule
Datos no aprendibles & (``unlearnable example*'' \textbf{OR} ``unlearnable data'' \textbf{OR} ``unlearnable dataset*'' \textbf{OR} ``unlearnable sample*'' \textbf{OR} ``unlearnability'' \textbf{OR} ``error-minimizing noise'' \textbf{OR} ``unexploitable data'' \textbf{OR} ``robust unlearnable example*'') \\
 & \textbf{OR} \\
Ataques de disponibilidad & ((``availability attack*'' \textbf{OR} ``availability poison*'') \textbf{AND} (image* \textbf{OR} ``training data'' \textbf{OR} ``deep learning'' \textbf{OR} ``machine learning'')) \\
 & \textbf{OR} \\
Perturbación protectora & (``data-protecting perturbation*'' \textbf{OR} ``protective perturbation*'' \textbf{OR} ``proactive perturbation*'') \\
 & \textbf{OR} \\
Envenenamiento indiscriminado & (``indiscriminate poison*'' \textbf{OR} ``contrastive poison*'') \\
 & \textbf{OR} \\
Camuflaje facial & (cloak* \textbf{AND} (``face recognition'' \textbf{OR} ``facial recognition'')) \\
\botrule
\end{tabularx}
\end{table}


\subsection{Criterios de selección}

Los criterios de inclusión (Tabla~\ref{tab:inclusion}) delimitan los estudios que pueden responder las preguntas de investigación. Los criterios de exclusión (Tabla~\ref{tab:exclusion}) descartan los estudios que, aun tratando el tema, no cumplen otras condiciones necesarias. Dos de ellos separan fenómenos próximos que la búsqueda recupera con frecuencia: las protecciones que actúan solo en la inferencia (CE6) y las dirigidas contra modelos generativos (CE7). El uso de un modelo generativo como contramedida para purificar los datos no activa el criterio CE7.

\begin{table}[t]
\caption{Criterios de inclusión}\label{tab:inclusion}
\footnotesize
\begin{tabularx}{\textwidth}{@{}L{0.9cm}Y{1.0}@{}}
\toprule
ID & Criterio \\
\midrule
CI1 & Estudio primario con evaluación empírica original, publicado como artículo completo en una revista o un congreso con revisión por pares. \\
CI2 & Publicado entre 2020 y 2026. \\
CI3 & Estudia datos visuales protegidos: imágenes, vídeos, pares imagen-texto con un componente visual pertinente o datos tridimensionales cuya tarea final sea discriminativa. \\
CI4 & Evalúa perturbaciones proactivas destinadas a dificultar el entrenamiento o ajuste no autorizado de modelos discriminativos, o contramedidas frente a ellas. \\
CI5 & Aporta evidencia pertinente para responder al menos una de las preguntas de investigación. \\
CI6 & El texto completo está en inglés o español. \\
\botrule
\end{tabularx}
\end{table}

\begin{table}[t]
\caption{Criterios de exclusión}\label{tab:exclusion}
\footnotesize
\begin{tabularx}{\textwidth}{@{}L{0.9cm}Y{1.0}@{}}
\toprule
ID & Criterio \\
\midrule
CE1 & No es una publicación primaria elegible: revisión, editorial, resumen sin artículo completo o preimpresión sin publicación arbitrada. \\
CE2 & Publicado fuera del periodo establecido. \\
CE3 & Texto completo en un idioma distinto del inglés o el español. \\
CE4 & No estudia datos visuales pertinentes al objeto de la revisión. \\
CE5 & Examina otro fenómeno, como el envenenamiento malicioso, las puertas traseras o el desaprendizaje de modelos, sin estudiar la protección proactiva. \\
CE6 & La protección actúa exclusivamente durante la inferencia de un modelo fijo, sin una evaluación pertinente al entrenamiento o ajuste. \\
CE7 & El objetivo es exclusivamente impedir el entrenamiento de modelos generativos. \\
CE8 & No presenta una evaluación empírica pertinente a ninguna pregunta de investigación. \\
\botrule
\end{tabularx}
\end{table}


La selección se realizó en dos etapas:
\begin{enumerate}
\item Tras eliminar los registros duplicados y las versiones de un mismo trabajo, dos revisores examinaron de forma independiente el título y el resumen de cada registro. Los registros sin una decisión clara pasaron a la lectura del texto completo.
\item En la etapa de elegibilidad, un revisor leyó el texto completo de cada estudio y el segundo revisor contrastó la decisión con la ficha de extracción y el acta de revisión del estudio.
\end{enumerate}
Las discrepancias se resolvían por consenso; no se presentaron discrepancias en ninguna de las dos etapas.

\subsection{Evaluación de la calidad}

La calidad de cada estudio incluido se valoró con las ocho preguntas de la Tabla~\ref{tab:calidad}. Las preguntas cubren la descripción técnica de la protección y de la contramedida, la identificación de datos y modelos, la separación entre desarrollo y prueba, la adecuación de los comparadores, la validez de las métricas, la incertidumbre, el reporte completo y la mesura de las conclusiones. La suma de las valoraciones, entre 0 y 8, no se utilizó para excluir estudios de la revisión, sino para decidir su participación en la síntesis.

\begin{table}[t]
\caption{Preguntas de evaluación de la calidad. Cada pregunta se valora como Sí (1), Parcial (0,5) o No (0); la información insuficiente y los casos no aplicables no puntúan}\label{tab:calidad}
\footnotesize
\begin{tabularx}{\textwidth}{@{}L{3.6cm}Y{1.0}@{}}
\toprule
Aspecto evaluado & Pregunta \\
\midrule
Descripción de la protección y la contramedida & ¿La protección y, cuando corresponde, la contramedida están descritas con su procedimiento y configuración técnica? \\
Identificación de datos y modelos & ¿Se identifican suficientemente los datos, la tarea, el modelo objetivo y el procedimiento de entrenamiento o ajuste? \\
Separación entre desarrollo y prueba & ¿La evaluación evita la contaminación indebida entre desarrollo y prueba, de acuerdo con la tarea y el escenario declarado? \\
Adecuación de los comparadores & ¿Los comparadores y las condiciones experimentales permiten interpretar el efecto evaluado bajo las capacidades declaradas del adversario? \\
Validez de las métricas & ¿Las métricas están definidas y miden el resultado que el estudio afirma demostrar? \\
Incertidumbre de los resultados & ¿Las repeticiones, las unidades evaluadas y las medidas de incertidumbre permiten valorar la estabilidad de los resultados? \\
Reporte completo & ¿Se presentan los resultados de las evaluaciones anunciadas y se justifican las omisiones o selecciones realizadas? \\
Mesura de las conclusiones & ¿Las conclusiones se limitan a lo demostrado y reconocen las limitaciones relevantes? \\
\botrule
\end{tabularx}
\end{table}


\subsection{Extracción de datos}

Los datos se extrajeron con un formulario común para todos los estudios (Tabla~\ref{tab:extraccion}). Cada valor se registró junto con su localización en la fuente (página, tabla o figura) y con su estado: informado, no informado, no evaluado por el estudio o no aplicable. La ausencia de una evaluación no se interpretó como un resultado nulo. Cuando un estudio presentaba varias configuraciones experimentales, se extrajeron por separado los resultados de cada una. \pend{Declarar el uso de herramientas de inteligencia artificial en la extracción y el control humano aplicado, conforme a la política editorial de la revista.}

\begin{table}[t]
\caption{Formulario de extracción de datos}\label{tab:extraccion}
\footnotesize
\begin{tabularx}{\textwidth}{@{}L{3.2cm}Y{1.0}@{}}
\toprule
Dimensión & Información extraída \\
\midrule
Protección & Nombre, mecanismo de generación, espacio y granularidad de la perturbación, dependencia de un modelo sustituto \\
Restricciones & Norma y presupuesto de la perturbación, proporción de datos protegidos \\
Escenario de evaluación & Modalidad, tarea, conjuntos de datos, arquitecturas y procedimiento de entrenamiento \\
Contramedidas & Nombre, familia, procedimiento, parámetros, etapa de aplicación y adaptación a la protección \\
Adversario & Conocimiento de la protección, acceso a datos, etiquetas y modelos, recursos y restricciones \\
Condiciones comparadas & Entrenamiento limpio, protegido, protegido con contramedida, limpio con la misma contramedida y recuperación autorizada \\
Medición & Métricas, dirección, unidad, umbral operativo y resultados por condición \\
Incertidumbre & Repeticiones, semillas y medidas de dispersión \\
Recuperación autorizada & Actor, recursos y resultado de la restauración legítima \\
Costos y efectos & Costo de generación y calidad visual del dato protegido \\
Conclusiones & Conclusiones del estudio y limitaciones declaradas \\
\botrule
\end{tabularx}
\end{table}


\subsection{Síntesis de los datos}\label{sec:sintesis}

\paragraph{Comparabilidad.} Cada par formado por la condición protegida y la condición con contramedida se clasificó según su comparabilidad:
\begin{itemize}
\item \emph{sustentada}, cuando ambas condiciones comparten conjunto de datos, arquitectura, protocolo de entrenamiento y métrica, y la contramedida está especificada;
\item \emph{limitada}, cuando falta uno de esos elementos;
\item \emph{no determinable}, cuando la fuente no permite establecerla.
\end{itemize}
La ausencia de repeticiones o de medidas de dispersión se consideró una limitación de la precisión del resultado, no de su comparabilidad. Las magnitudes de la Sección~\ref{sec:condiciones} se calcularon solo para los pares con comparabilidad sustentada.

\paragraph{Niveles de síntesis.} Todos los estudios incluidos se caracterizaron para responder RQ1, RQ2 y RQ4. Participaron en la síntesis los estudios con una valoración de calidad igual o superior a 4,5 sobre 8 o con al menos una comparación sustentada. La síntesis cuantitativa de RQ3 se basó en los estudios con al menos una comparación sustentada.

\paragraph{Agrupación.} Las protecciones se agruparon según su estrategia de generación y las contramedidas según su familia. Cuando una condición combinaba contramedidas de varias familias, se asignó a la familia con mayor capacidad esperada de eliminar la perturbación, en este orden: purificación generativa, aumentos adversarios, entrenamiento adversario, transformación y compresión, aumentos de datos, proyección y sondas lineales, cambio de paradigma y detección.

\paragraph{Agregación.} Las comparaciones de un mismo estudio no son independientes y su número varía mucho entre estudios. Por ello, la agregación se realizó en dos pasos: primero se obtuvo la mediana dentro de cada estudio y después la mediana y el rango intercuartílico entre estudios. Como indicador de neutralización se calculó, para cada estudio, la proporción de comparaciones que quedan a 5~pp o menos del entrenamiento limpio, y se promedió entre estudios; los umbrales de 2 y 10~pp se informan como análisis de sensibilidad. No se realizó un metaanálisis, debido a la dependencia entre comparaciones, a la escasez de medidas de dispersión y a la heterogeneidad de las configuraciones.
